\documentclass[UTF8,12pt,a4paper]{ctexart}

\usepackage[a4paper,margin=2.5cm]{geometry}
\usepackage{longtable,array,xcolor,hyperref}
\hypersetup{hidelinks}
\setlength{\parindent}{2em}
\pagestyle{plain}
\newcommand{\TODO}[1]{\textcolor{red}{\textbf{TODO[#1]}}}
\newcolumntype{L}[1]{>{\raggedright\arraybackslash}p{#1}}

\begin{document}
\begin{center}
  {\zihao{2}\bfseries AI工具使用详情}
\end{center}

\section{工具与使用范围}
\TODO{逐项填写实际工具名称、版本、模型、提供方，以及使用阶段和用途。}

\section{主要提示方式与过程}
\TODO{概述如何提供题目、代码和证据约束以及实际使用过程；不得虚构未发生的交互。必要时可附典型交互示例。}

\section{采纳、人工修改与核验}
\renewcommand{\arraystretch}{1.3}
\begin{longtable}{@{}L{0.06\textwidth}L{0.13\textwidth}L{0.15\textwidth}L{0.23\textwidth}L{0.25\textwidth}@{}}
  \hline
  序号\newline 时间 & 工具、版本\newline 与模型 & 目的\newline 与阶段 & 采纳内容及人工修改 & 核验方式与结果 \\
  \hline
  1 & \TODO{填写} & \TODO{填写} & \TODO{填写} & \TODO{填写} \\
  \hline
\end{longtable}

\section{责任确认}
\TODO{如实确认AI参与深度、人工复核范围及论文/代码/数据一致性；不得把AI完成的核心建模改写为团队独立完成。}
\end{document}
